\section{Conclusion}\label{sec:conclusion}

This paper has asked when a fiat-backed stablecoin behaves as the par-convertible payment instrument it promises to be, when it does not, and through which structural channel the answer changes. We build a microstructure model in which heterogeneous user types face a Morris-Shin coordination problem at the conversion stage and choose between converting at the issuer and selling on an exchange. The strategic stage is solved as a fixed point on the type-specific thresholds. We calibrate the model to the March 2023 USDC episode at hourly frequency, evaluate the calibrated structural block out of sample on USDT and DAI, falsify it on TerraUSD, and map it into a per-corridor cost comparison against incumbent payment services.

\paragraph{When the stablecoin behaves as expected.} The model rationalises a tranquil baseline regime as a high-precision prior on the issuer fundamental: when uncertainty about reserves is low, the calibrated parameters produce a secondary-market price that sits within tens of basis points of par across the entire Monte Carlo sample. Empirically this corresponds to about ninety-nine percent of hours in the 2022 calibration window. In the per-corridor cost-benefit comparison, this is the regime in which the fee savings from using a stablecoin dominate the priced risk. The retail and cross-border corridors deliver a negative net metric of $339$ and $294$ basis points respectively even after probability-weighting the regimes by their empirical frequency.

\paragraph{When the stablecoin does not behave as expected, and through which channel.} The same calibrated parameters produce a stress regime under a low-precision prior on the fundamental, with prices on the exchange that fan out into a fat-tailed distribution. Three structural channels carry the cost. The expected discount conditional on choosing to convert rises sharply with the stress-regime $\theta$-distribution, contributing roughly $200$ basis points of additional cost on average. The default-loss channel scales with the issuer hazard, contributing about $380$ basis points more. The congestion channel does not change with the regime but contributes the most in the wholesale corridor through the type-specific congestion-cost coefficient. Proposition~\ref{prop:reserve} characterises the secondary-market channel formally: tighter reserves dampen aggregate conversion pressure at the fixed point because the anticipated price discount makes conversion less attractive at the margin. Proposition~\ref{prop:convenience} characterises the payment-use channel in the opposite direction: more useful stablecoins are more robust at the conversion stage. The two propositions deliver clean theoretical statements of the empirical mechanism documented by \citet{ma_zeng_zhang_2025} and \citet{bertsch_2025_adoption_fragility}.

\paragraph{The boundary of the model and what it identifies.} The calibrated USDC structural block transports to DAI with a tight fit on the average severity of the stress regime, consistent with DAI's exposure to USDC through its collateral. It does not transport to USDT, whose par-stable trajectory in the March 2023 window is not reproducible at the calibrated block. The mismatch is informative as a boundary statement: the framework describes stablecoins with open primary-market arbitrage and reserve-backed conversion, and does not describe stablecoins with concentrated arbitrage access or with algorithmic backing. The boundary is precisely the structural difference that \citet{ma_zeng_zhang_2025} document for USDT and that algorithmic stablecoins display in their pre-collapse trajectory. The TerraUSD falsification confirms the same boundary at the level of an explicit event: with the issuer reserve set to zero, the calibrated model cannot reproduce the realised collapse, in which the trajectory falls to roughly five cents per unit, because the model's redemption stage does not represent the recursive feedback through the paired backing token.

\paragraph{Comparing stablecoins to incumbent services dynamically.} The cost-benefit metric reported in Section~\ref{sec:cost_benefit} is a per-period static building block, not a dynamic adoption path. The static piece is in place. Embedding $\mathrm{CEC}^{stbl}_c(t)$ and $\mathrm{CEC}^{inc}_c(t)$ in a discrete choice over payment services at the user level gives a logit-share equation $\sigma_c(t) = \exp(-\lambda \mathrm{CEC}^{stbl}_c(t)) / (\exp(-\lambda \mathrm{CEC}^{stbl}_c(t)) + \exp(-\lambda \mathrm{CEC}^{inc}_c(t)))$, with $\lambda$ a user-level sensitivity to the cost differential that can be estimated from transaction-level data on the marginal user's choice in the corridor. Network effects on the stablecoin side enter through the dependence of the secondary-market depth $M$ on the adoption share, which feeds back into the strategic stage and, through Proposition~\ref{prop:reserve}, into the implicit reserve adequacy of the issuer. Iterating the share equation forward with the depth feedback would identify two policy-relevant objects that the static metric cannot. The first is the network threshold at which a corridor becomes self-sustaining for the stablecoin, in the sense that an additional marginal user finds the stablecoin attractive precisely because the existing adoption raises $M$ enough to lower the priced tail risk. The second is the fitted adoption curve for a corridor with an exogenous cost shock, against which the realised post-package flow displacement after 2 February 2026 in Brazil provides an observable benchmark once a longer post-package window is available. We provide the static metric and the calibrated structural parameters; the dynamic extension is a natural follow-on.

\paragraph{Policy takeaways.} The calibrated framework supports three takeaways at the policy level. The first is that the cost comparison against an incumbent payment service is corridor-specific and regime-conditional, and authorities or institutional users should not draw a universal verdict on stablecoins from one corridor's number. In the calibrated environment, the retail and cross-border net metrics are favourable to the stablecoin by margins that survive a tenfold increase in stress frequency; the wholesale net metric is favourable to the incumbent under settlement-finality assumptions that wholesale users themselves must verify in context. The second is that the framework is reusable. Plugging new fee data, new corridor definitions, or new regulatory cost wedges into the calibrated module updates the per-corridor metric without re-estimating the structural block, as long as the corridor's issuer falls within the open-arbitrage stablecoin segment that the model describes. The third is that the framework allows discrete policy events to be evaluated as marginal cost shocks, and Section~\ref{sec:cost_benefit} reports the magnitude of the friction required to equalise the two costs in the cross-border corridor under the November 2025 Brazilian package. A payment-policy authority calibrating the strength of a similar package can read the sensitivity table in Appendix~\ref{sec:sensitivity} as a basis-point-by-basis-point map between the chosen friction and the resulting equalisation point.

\paragraph{What is left.} Three extensions deserve attention. The first is the structural extension to stablecoins with concentrated primary-market arbitrage, which would generalise the calibrated microstructure block to cover USDT and would close the OOS gap. The second is the dynamic extension just described, which would convert the static cost-benefit metric into an explicit path of adoption shares per corridor and would close the loop with the November 2025 Brazilian flow data once a longer post-package window is available. The third is the algorithmic extension that would add the recursive backing-token feedback documented in the TerraUSD falsification and would let the model identify the conditions under which algorithmic backing transitions from a stable equilibrium to a self-reinforcing collapse.
